\chapter{Preliminares}
\label{chap:preliminares}

% Introduccion cohesiva temporal

En este capítulo se presentan los conceptos fundamentales de la teoría de lenguajes formales que serán necesarios a lo largo de este trabajo. Se introducen los elementos básicos como alfabetos y cadenas, las operaciones sobre cadenas y lenguajes, así como los problemas clásicos que se estudian en este ámbito.

\section{Teoría de lenguajes}

% Introduccion cohesiva temporal

La teoría de lenguajes formales estudia los conjuntos de cadenas construidas a partir de un alfabeto y las operaciones que se pueden realizar sobre ellas. Esta teoría constituye la base matemática de áreas como la computación, la compilación y la lógica formal.

\subsection{Conceptos, operaciones y problemas relacionados con lenguajes}

El concepto más fundamental en teoría de lenguajes es el de alfabeto, que sirve como conjunto de construcción para las cadenas.

\begin{definition}{(Alfabeto)}
	Un alfabeto, usualmente denotado como $\alphabet$, es un conjunto finito no vacío de elementos llamados símbolos.
\end{definition}

A partir de un alfabeto, se puede definir una \emph{cadena}, que corresponde a una secuencia finita de símbolos de dicho alfabeto.

\begin{definition}{(Cadena)}
	Una cadena sobre un alfabeto es una sucesión finita de símbolos pertenecientes a ese alfabeto.
\end{definition}

Existe una cadena especial denominada \emph{cadena vacía}, que no contiene ningún símbolo.

\begin{definition}{(Cadena vacía)}
	La cadena vacía es aquella que no posee ningún símbolo y se denota por $\eword$.
\end{definition}

Se denota por $|w|$ a la longitud de una cadena $w$. La cadena vacía $\eword$ tiene longitud cero.

La operación principal sobre cadenas es la concatenación, que permite unir dos cadenas de manera secuencial.

\begin{definition}{(Concatenación)}
	Sean $x$ y $y$ dos cadenas. La concatenación de $x$ y $y$, denotada por $xy$, es la cadena formada por todos los símbolos de $x$ seguidos de todos los símbolos de $y$.
\end{definition}

Si se aplica la concatenación de manera repetida sobre una misma cadena, se obtiene lo que se conoce como \emph{potencia} de la cadena.

\begin{definition}{(Copia o potencia)}
	Dada una cadena $w$ y un entero $n \geq 0$, la potencia de $w$ se define inductivamente como:

	\textbf{Caso base:} $w^0 = \eword$

	\textbf{Paso inductivo:} $w^n = w w^{\,n-1}$ para $n > 0$.
\end{definition}

Con estos conceptos se puede definir formalmente un lenguaje, que no es más que un conjunto de cadenas.

\begin{definition}{(Lenguaje)}
	Dado un alfabeto $\alphabet$, un lenguaje sobre $\alphabet$ es cualquier conjunto de cadenas formadas con símbolos de $\alphabet$.
\end{definition}

Denotaremos por $\allwords$ al lenguaje de todas las cadenas posibles sobre un alfabeto dado, y por $\allnewords$ al lenguaje de todas las cadenas no vacías, es decir, $\allnewords = \allwords - \{\eword\}$.

\paragraph{Ejemplo} Sea el alfabeto $\alphabet = \{0, 1\}$. Algunas cadenas sobre este alfabeto son $0$, $1$, $10$, $1101$ y $\eword$. Un ejemplo de lenguaje sobre $\alphabet$ es el conjunto de todas las cadenas que terminan en 0.

Al ser los lenguajes conjuntos de cadenas, es posible definir sobre ellos operaciones clásicas de conjuntos como unión, intersección y complemento \cite{automataTheory}. Además, resulta interesante estudiar ciertas propiedades y problemas asociados a los lenguajes, tales como verificar si una cadena pertenece a un lenguaje o si un lenguaje es vacío.

\begin{definition}{(Problema del vacío)}
	El problema del vacío consiste en determinar, para un lenguaje dado $L$, si $L = \emptyset$.
\end{definition}

\subsection{Gramáticas y Jerarquía de Chomsky}
\label{sec:grammars}

Existen diferentes formas de representar un lenguaje, entre ellas las grámaticas constituyen un formalismo para representar lenguajes, de forma concisa y clara.

\begin{definition}{(Gramática)}
	Una gramática es un formalismo definido mediante una 4-tupla:
	$$
		G = (N, T, P, S)
	$$

	donde:
	\begin{itemize}
		\item $N$ es un conjunto de simbolos no terminales.
		\item $T$ es un conjunto de símbolos terminales. Debe cumplirse que $N \cup T = \emptyset$
		\item $P$ es un conjunto de pares ordenados de la forma:
		      $$
			      (\alpha, \beta)
		      $$
		      tales que $\alpha \in (N \cup T)^* - T^*$ y $\beta \in (N \cup T)^*$ llamados reglas de producción.

		      Las reglas $(\alpha, \beta)$ se denotan como $\alpha \gprod \beta$
		\item $S$ es el símbolo inicial de la gramática, $S \in N$

	\end{itemize}
\end{definition}

Para que la definición anterior sea útil para generar lenguajes es necesario definir el proceso con el que se generan cadenas a partir de una gramática.

\begin{definition}{(Derivación directa)}
	Sean $x,y \in (N \cup T)^*$ existe una derivacíon de $x$ en $y$ si $x = x_1 \alpha x_2$ y $y = y_1 \beta y_2$ y en $G$ existe la regla $\alpha \gprod \beta$. Esto se denota como $x \gder y$
\end{definition}

Una derivacíon no es mas que una sustitición de una subcadena, de la cadena izquierda por otra, dando lugar a la cadena de la derecha. Este proceso se puede generalizar a mas de un paso.

\begin{definition}{(Derivación)}
	Sean $x,y \in (N \cup T)^*$ existe una derivacíon de $x$ en $y$ si existen $y_1, \dots, y_n$ tales que:
	$$
		x \gder y_1, \gder \dots \gder y_n = y
	$$

	lo cual se denota por $x \gderp y$ y por $x \gderz y$ para incluir el caso de cero pasos de derivación.
\end{definition}

Con esto es posible definir el lenguaje generado por una gramática.

\begin{definition}{(Lenguaje generado)}
	Sea $G$ una gramatica, el lenguaje generado por $G$ es el siguiente conjunto:
	$$
		L(G) = \{x \suchthat x \in T \land S \gderp x\}
	$$
\end{definition}

No todas las gramáticas o formalismos para generar lenguajes son iguales. La jerarquía de Chomsky \cite{hunter2020chomsky} clasifica las gramáticas en diferentes tipos según la forma de sus reglas de producción, lo que a su vez determina el tipo de lenguajes que pueden generar. Estas gramaticas se incluyen de forma estricta, es decir, cada tipo de gramática es un subconjunto estricto del tipo anterior:

\begin{itemize}
	\item Las gramáticas \textbf{tipo-0}, también conocidas como gramáticas irrestrictas, son las más generales. Sus reglas no tienen restricciones (son de la forma $\alpha \gprod \beta$ donde $\alpha \in (N \cup T)^* - T^*$ y $\beta \in (N \cup T)^*$) y generan lenguajes recursivamente enumerables (\textit{RE}).

	\item En las gramáticas de \textbf{tipo-1}, o gramáticas sensibles al contexto \text{CS}, las reglas de producción son de la forma $\alpha A \gamma \gprod \alpha \beta \gamma$, donde $A \in N$, $\alpha, \beta, \gamma \in (N \cup T)^*$, y \(|\beta| \geq 1\). Estas generan los lenguajes dependientes del context.

	\item Las gramáticas de \textbf{tipo-2}, o gramáticas libres de contexto (\textit{CFG}), en la construcción de compiladores. Sus reglas de producción son de la forma $A \gprod \beta$, donde $A \in N$ y $\beta \in (N \cup T)^*$. Estas generan los lenguajes libres de contexto, que son reconocidos por autómatas de pila \cite{automataTheory}.

	\item Finalmente, las gramáticas de \textbf{tipo-3}, o gramáticas regulares, son con reglas de producción de la forma $A \gprod aB$ o $A \gprod a$, donde $A, B \in N$ y $a \in T$. Estas gramáticas generan los lenguajes regulares (\textit{REG}), que son reconocidos por autómatas finitos \cite{automataTheory}.
\end{itemize}

Un formalismo o gramatica es más poderoso que otro si puede generar un conjunto de lenguajes que el otro no puede, o sea, la familia de lenguajes generados es mas grande.

\begin{definition}{(Familia de lenguajes)}
	Sea $X \in \{\re, \csen, \cfg, \reg\}$ una clase de gramáticas, la familia de lenguajes generados por $X$ se define como:
	$$
		\mathcal{L}(X) = \{L(G) \suchthat G \text{ es una gramática de tipo } X\}
	$$
\end{definition}

La jerarquía de Chomsky establece que:
$$
	\mathcal{L}(\reg) \subset \mathcal{L}(\cfg) \subset \mathcal{L}(\csen) \subset \mathcal{L}(\re)
$$

Como ejemplos de la inclusion estricta de estas familias de lenguajes, se pueden mencionar los siguientes:

\begin{itemize}
	\item El lenguaje $L_1 = \{a^n b^n \suchthat n \geq 0\}$ es un lenguaje libre de contexto que no es regular, por lo que $L_1 \in \mathcal{L}(\cfg) - \mathcal{L}(\reg)$ \cite{automataTheory}.

	\item El lenguaje $L_2 = \{ww \suchthat w \in \{a, b\}^*\}$ es un lenguaje sensible al contexto que no es libre de contexto, por lo que $L_2 \in \mathcal{L}(\csen) - \mathcal{L}(\cfg)$ \cite{automataTheory}.
\end{itemize}

\subsection{Lenguajes regulares y autómatas}
\label{sec:automata}

Un autómata es un modelo matemático para describir sistemas que pueden estar en diferentes estados y cambiar de estado en respuesta a entradas. En la teoría de lenguajes formales, los autómatas se utilizan para determinar si una cadena pertenece a un lenguaje específico. Cada grámatica de la jerarquía de Chomsky tiene un tipo de autómata asociado que reconoce exactamente los lenguajes generados por esa gramática. A continuación se introduce el autómata finito determinista, que es el modelo de autómata asociado a los lenguajes regulares.

\begin{definition}{(Autómata finito determinista)}
	Un autómata finito determinista (AFD) \cite{automataTheory} es una 5-tupla:
	$$
		\mathcal{A} = (Q, \alphabet, \delta, q_0, F)
	$$

	donde:
	\begin{itemize}
		\item $Q$ es un conjunto finito de estados.
		\item $\alphabet$ es un alfabeto de entrada.
		\item $\delta: Q \times \alphabet \to Q$ una función de transición que define el cambio de estado para cada símbolo de entrada.
		\item $q_0 \in Q$ es el estado inicial.
		\item $F \subseteq Q$ es el conjunto de estados de aceptación o finales.
	\end{itemize}
\end{definition}

Un AFD procesa una cadena de entrada comenzando en el estado inicial $q_0$ y aplicando la función de transición $\delta$ para cada símbolo de la cadena. Si al finalizar el procesamiento de la cadena el autómata se encuentra en un estado que pertenece a $F$, entonces se dice que el autómata acepta la cadena, lo que implica que la cadena pertenece al lenguaje reconocido por el autómata. Para formalizar esto, se define la función de transición extendida $\hat{\delta}$ que permite procesar cadenas completas.

\begin{definition}{(Función de transición extendida)}
	La función de transición extendida $\hat{\delta}: Q \times \alphabet^* \to Q$ se define inductivamente como sigue:

	\textbf{Caso base:} Para cualquier estado $q \in Q$, $\hat{\delta}(q, \eword) = q$.

	\textbf{Paso inductivo:} Para cualquier estado $q \in Q$, símbolo $a \in \alphabet$ y cadena $w \in \alphabet^*$, se define:
	$$
		\hat{\delta}(q, aw) = \hat{\delta}(\delta(q, a), w)
	$$
\end{definition}

Con esta función de transición extendida, se define ellenguaje reconocido por un AFD.

\begin{definition}{(Lenguaje reconocido por un AFD)}
	El lenguaje reconocido por un AFD $\mathcal{A}$ es el conjunto de todas las cadenas que el autómata acepta, es decir:
	$$
		L(\mathcal{A}) = \{w \in \alphabet^* \suchthat \hat{\delta}(q_0, w) \in F\}
	$$
\end{definition}

En \cite{automataTheory} se demuestra que los lenguajes reconocidos por AFDs son exactamente los lenguajes regulares, es decir, $L(\mathcal{A}) \in \mathcal{L}(\reg)$ para cualquier AFD $\mathcal{A}$, y para cada lenguaje regular $L$ existe un AFD $\mathcal{A}$ tal que $L = L(\mathcal{A})$.

% TODO: mejorar esta introduccion temporal cuando tenga el capitulo
Los lenguajes regulares y autómatas finitos deterministas seran importantes a lo largo de este trabajo, ya que nos permitirán modelar el lenguaje las interpretaciones de formulas booleanas verdaderas sin restricciones en las variables luego.

\section{Complejidad computacional}
La teoría de la complejidad computacional se ocupa de clasificar los problemas según la cantidad de recursos computacionales que requieren para ser resueltos.

\subsection{Notación asintótica}

La notación asintótica se usa para describir el comportamiento de una función $f(n)$, donde $n$ es el tamaño de la entrada. A continuación se define la notación $O$ (grande):

\begin{definition}{(Notación $O$)}
	Sea $f(n)$ y $g(n)$ dos funciones de $n$. Se dice que $f(n) = O(g(n))$ si existen constantes positivas $c$ y $n_0$ tales que:

	$$
		f(n) \leq c \cdot g(n) \quad \text{para todo $n \geq n_0$}
	$$
\end{definition}

Esta notación se utiliza para describir el límite superior del crecimiento de una función, lo que es útil para analizar la eficiencia de algoritmos y la complejidad de problemas computacionales.

Se dice que un algoritmo requiere una cantidad de tiempo polinomial si su tiempo de ejecución es $O(n^k)$ para alguna constante $k$ y entrada de tamaño $n$. La complejidad de un algoritmo es la complejidad del algoritmo mas eficiente que lo resuelve en el peor caso.

\subsection{Clases de problemas}

Los problemas de decisión se agrupan en diferentes clases segun los recursos computacionales necesarios para resolverlos. A continuación se definen las clases de complejidad que se utilizarán a lo largo de este trabajo:

\begin{definition}{(Clase P)}
	La clase de complejidad P es el conjunto de problemas de decisión que pueden ser resueltos por un algoritmo determinista en tiempo polinomial \cite{automataTheory}.
\end{definition}

\begin{definition}{(Clase NP)}
	La clase de complejidad NP es el conjunto de problemas de decisión para los cuales una solución propuesta puede ser verificada por un algoritmo determinista en tiempo polinomial \cite{automataTheory}.
\end{definition}

\begin{definition}({Clase NP-duro})
	Un problema de decisión es NP-duro si cada problema en NP puede ser reducido a él en tiempo polinomial \cite{automataTheory}.
\end{definition}

\begin{definition}({Clase NP-completo})
	Un problema de decisión es NP-completo si pertenece a NP y es NP-duro \cite{automataTheory}.
\end{definition}

Posiblemente sea el problema abierto más famoso en las ciencias de la computación determinar si P es igual a NP. Es conocido que P es un subconjunto de NP, pero hasta la fecha no se sabe si esta inclusión es estricta o si ambos conjuntos son iguales \cite{automataTheory}. El conjunto de los problemas NP-completos es de particular interés, ya que si se encuentra un algoritmo de tiempo polinomial para resolver cualquier problema NP-completo, entonces se demostraría que P es igual a NP. Por otro lado, si se demuestra que no existe tal algoritmo, entonces se confirmaría que P es diferente de NP.

A continuación se presenta el problema de satisfacibilidad booleana, que es un problema NP-completo y de gran relevancia para este trabajo.


\section{El problema de la satisfacibilidad booleana}

El problema de la satisfacibilidad booleana (SAT) es un problema de decisión que consiste en determinar si existe una asignación de valores de verdad a las variables de una fórmula booleana que haga que la fórmula sea verdadera \cite{automataTheory}.

\subsection{Definiciones}

A continuación se definen los conceptos necesarios para formalizar el problema de satisfacibilidad booleana.

\begin{definition}{(Variable booleana)}
	Una variable booleana es una variable que puede tomar solo dos valores: $1$ (verdadero) o $0$ (falso).
\end{definition}

\begin{definition}{(Literal)}
	Sea $x$ una variable booleana, un literal o término es una variable o su negación. Es decir, $x$ o $\neg x$ son literales.
\end{definition}

\begin{definition}{(Cláusula)}
	Una cláusula $C$ es una disyunción finita de literales, es decir:
	$$
		t_1 \lor t_2 \lor \dots \lor t_k
	$$
\end{definition}

\begin{definition}{(Fórmula en forma normal conjuntiva)}
	Una fórmula booleana está en forma normal conjuntiva (CNF) si es una conjunción de cláusulas. Es decir, una fórmula $F$ está en CNF si se puede escribir como:
	$$
		F = C_1 \land C_2 \land \dots \land C_m
	$$
	donde cada $C_i$ es una cláusula.

\end{definition}

En lo que sigue, se considerarán únicamente fórmulas booleanas en forma normal conjuntiva. Por simplicidad, se las denominará simplemente fórmulas booleanas o fórmulas.


\begin{definition}{(Interpretación)}
	Una interpretación es una asignación de valores de verdad a las variables de una fórmula booleana. Formalmente, sea $F$ una fórmula y sea $X = \{x_1, \ldots, x_n\}$ el conjunto de variables de $F$, una interpretación es una función $v: X \to \{0, 1\}$.
\end{definition}

\begin{definition}{(Evaluación de una fórmula)}
	Sea $v: X \to \{0,1\}$ una interpretación. La evaluación de una fórmula booleana bajo $v$ es la función
	$$
		(\cdot)_v : \mathcal{F} \to \{0,1\}
	$$

	donde $\mathcal{F}$ es el conjunto de todas las fórmulas booleanas, y se
	define como sigue:

	\begin{itemize}
		\item Si $x$ es una variable, entonces $(x)_v = v(x)$ y
		      $$
			      (\neg x)_v =
			      \begin{cases}
				      1 & \text{si } v(x)= 0,  \\
				      0 & \text{si } v(x) = 1.
			      \end{cases}
		      $$

		\item Una cláusula es verdadera si al menos uno de sus literales es verdadero. Formalemente, si $C = t_1 \lor \dots \lor t_k$, entonces
		      $$
			      (C)_v = 1 \iff \exists i \in \{1,\dots,k\} \text{ tal que } (t_i)_v = 1.
		      $$

		\item Una fórmula es verdadera si todas sus cláusulas son verdaderas. Formalemente, si $F = C_1 \land \dots \land C_m$, entonces
		      $$
			      (F)_v = 1 \iff \forall i \in \{1,\dots,m\},\ (C_i)_v = 1.
		      $$
	\end{itemize}
\end{definition}


\begin{definition}{(Problema de satisfacibilidad booleana)}
	Sea $F$ una fórmula booleana en CNF, el problema de satisfacibilidad booleana (SAT) consiste en determinar si existe una interpretación $v$ para las variables de $F$ tal que $(F)_v = 1$.
\end{definition}

Este problema servirá como base para los resultados que se presentarán en este trabajo.

\subsection{SAT como problema NP-completo}

El problema de satisfacibilidad booleana es uno de los problemas más relevantes en la teoría de la complejidad computacional. En particular, fue el primer problema demostrado como NP-completo por Cook en \cite{Cook1971}, lo que lo convierte en un punto de referencia fundamental para el estudio de la complejidad computacional.

La importancia de SAT radica en que cualquier problema en NP puede transformarse, en tiempo polinomial, en una instancia equivalente de satisfacibilidad booleana \cite{automataTheory}. En consecuencia, SAT actúa como un problema representativo de la clase NP: resolverlo eficientemente implicaría la existencia de algoritmos eficientes para todos los problemas en NP.

% TODO: hablar de 2-sat y otras variantes polinomiales

\subsection{SAT desde la perspectiva de lenguajes formales}

El problema SAT puede ser abordado desde la perspectiva de lenguajes formales. En este sentido en la Facultad de Matemática y Computación de la Universidad de La Habana se han realizado varios trabajos: \cite{aCFSAT}, \cite{aSRCSAT}, \cite{MSAT} y \cite{RCGSAT}, los cuales modelan el problema como un lenguaje formal y buscan resolver instancias específicas en tiempo polinomial o demostrar que ciertas clases de lenguajes tienen problemas de decisión NP-duros.

En estos trabajos se toman dos enfoques principales utilizando el problema del vacío o de la palabra en lenguajes respectivamente. En el primer caso, se construye el lenguaje de las interpretaciones verdaras de una fórmula y se determina si este es vacío o no, en caso de no ser vacío implica que existe una interpretación verdadera y que la formula es satisfacible. En el segundo caso se construye el lenguaje de las fórmulas satisfacibles y se determina si la fórmula dada pertenece a ese lenguaje.

En \cite{aCFSAT} primero se asume que todas las variables de la formula son distintas y se construye el lenguaje de las interpretaciones verdaderas de la fórmula, sin dichas restricciones en las variables. Se demuestra que dicho lenguaje es regular mediante la construcción de un autómata finito, el cuál se denominó autómata booleano, que representa las reglas de la lógica proposicional en sus transiciones y sus estados representan un valor de verdad para la fórmula que se tiene hasta el momento. Por último se intersecta dicho automata con un lenguaje libre del contexto que garantize que todas las instancias de la misma variable tengan el mismo valor de verdad, y se determina si el lenguaje resultante es vacío o no.

En \cite{aSRCSAT} se intersecta el autómata booleano con grámaticas de concatenacion de rango simple lo cual amplia el conjunto de instancias se SAT solubles en tiempo polinomial. En \cite{MSAT} se aplica una idea similar, intersectando el autómata booleano con gramáticas matriciales simples, que aunque en el caso general son de tamaño exponencial con respecto a la fórmula, introduce la idea de utilizar una variante del lenguaje \textit{Copy} para garantizar que todas las instancias de la misma variable tengan el mismo valor de verdad.

En \cite{RCGSAT} se toma el enfoque del problema de la palabra. Se construye el lenguaje de las fórmulas satisfacibles mediante una transducción de una variante del lenguaje \textit{Copy}, y además se propone una gramática de concatenación de rango que genera reconoce este lenguaje. Como resultado se demuestra que el problema de la palabra para todos los formalismos que generen la variante del lenguaje \textit{Copy} y sean cerrados bajo transducción finita es Np-duro.

% TODO: leer articulos otros articulos y mencionarlos aquí